\section{Mode D2 -- Diagnostic / traitement de la défaillance}
\label{secModeD2}

\subsection{Conditions d'entrée}

Le mode D2 est atteint depuis le mode D1, dès que la communication avec
l'ensemble des équipements de la cellule est redevenue opérationnelle.

\subsection{Comportement attendu}

\begin{itemize}
    \item \textbf{Robot} : l'installation doit amener le robot dans une
    configuration saine, exempte de tout défaut, sans mouvement en attente
    d'exécution, prête à recevoir de nouvelles commandes, et avec le bras
    remis sous puissance.
    \item \textbf{Convoyeur / guichet} : l'accès au guichet reste interdit,
    comme en mode D1. Le système de préhension conserve son état.
    \item \textbf{IHM} : la situation de diagnostic/traitement de la
    défaillance doit être signalée à l'opérateur par une signalisation
    dédiée, clignotante, afin de la distinguer de l'arrêt d'urgence fixe
    du mode D1.
\end{itemize}

\begin{UPSTIidee}{Mode transitoire}
Ce mode a vocation à être transitoire : il ne fait qu'assainir la
configuration de l'installation avant de permettre la reprise de
l'activité. Aucune action de production n'y est réalisée.
\end{UPSTIidee}

\subsection{Conditions de sortie}

\begin{itemize}
    \item Vers le mode \textbf{A5} (\emph{Préparation pour remise en route
    après défaillance}), une fois la configuration saine de l'installation
    obtenue.
\end{itemize}
